\chapter{负载 关节零位 工具与相机标定}\label{ch:calibration}
\section{先明确究竟是哪一个模型不准确}
机械臂末端偏离目标，不一定是同一种误差。关节反馈角与驱动器目标角不一致，属于跟踪或执行问题；角度一致而末端实测位置不一致，可能来自几何、零位、工具变换或结构弹性；位置基本正确而低阻抗时向某方向偏，可能与重力、负载、摩擦或外力估计有关。

标定（calibration）是在一个选定模型内，通过参考测量确定参数。它与闭环控制的职责不同：控制器尽力实现自己知道的目标，标定则让“控制器知道的位置或力”与真实物理量一致。没有独立参考，只比较同一套模型输出与该模型自己的目标，不能证明绝对精度。

\begin{longtable}{p{0.20\linewidth}p{0.32\linewidth}p{0.38\linewidth}}
\toprule 问题&主要未知参数&需要的参考信息\\\midrule\endhead
负载重力补偿&负载质量及一阶质量矩&不同姿态的力矩或六维力传感器数据\\
负载动态补偿&负载质量 质心与惯量&足够变化的平动和转动加速度\\
关节零位标定&编码器零点与实际几何零点差&外部位置或姿态测量\\
TCP 标定&法兰到工具工作点的位置&同一固定点在不同法兰姿态下的约束\\
手眼标定&法兰与相机之间的刚体变换&机器人相对运动与相机相对运动\\
相机内参标定&焦距 主点及畸变参数&图像中的已知几何目标\\\bottomrule
\end{longtable}

\section{负载重力为什么只需要四个惯性组合}
\subsection{先把重力写在法兰坐标系}
设世界或基座坐标系为 $B$，法兰坐标系为 $F$。$R_{BF}$ 把法兰分量变成基座分量，重力加速度在基座中为 $a_g^B$。则
\[
 a_g^F=R_{BF}^{\mathsf T}a_g^B.
\]
负载质量为 $m_p$，质心在法兰坐标中的位置为 $c_p$，一阶质量矩为 $h_p=m_pc_p$。负载重力作用到法兰的力旋量，按本书 $[f;m]$ 顺序写为
\begin{equation}
 w_g^F=\begin{bmatrix}m_pa_g^F\\h_p\times a_g^F\end{bmatrix}
 =\underbrace{\begin{bmatrix}a_g^F&0_{3\times3}\\0_{3\times1}&-[a_g^F]_\times\end{bmatrix}}_{A_p(R_{BF})}
 \begin{bmatrix}m_p\\h_p\end{bmatrix}.
\end{equation}
$[a]_\times b=a\times b$。因为 $h\times a=-a\times h$，矩阵下半部分出现负号。四个未知量就是质量一个数、一阶矩三个数；静止数据不包含角加速度和角速度，无法据此确定全部转动惯量。

如果 $J_F$ 将关节速度映射到关于法兰点、在法兰系表达的速度旋量，负载重力对机械臂的广义作用力是 $J_F^{\mathsf T}w_g^F$，所需重力补偿为
\begin{equation}
 \Delta g(q)=-J_F(q)^{\mathsf T}w_g^F.
\end{equation}
这是“环境重力作用力”与“抵消它的驱动力矩”的区别。测力传感器给出的究竟是哪一侧作用哪一侧的力，要按其符号约定核对，不能把传感器读数直接当成补偿力矩。

\begin{examplebox}[算例 10A 两种方向就能看出一阶质量矩的作用]
假设负载 $m_p=2\ \mathrm{kg}$，$h_p=(0.2,-0.1,0.1)^{\mathsf T}\ \mathrm{kg\,m}$。第一姿态的法兰系重力为 $(0,0,-9.81)^{\mathsf T}$，则
\[
 f_1=(0,0,-19.62)^{\mathsf T}\ \mathrm N,\qquad
 m_1=(0.981,1.962,0)^{\mathsf T}\ \mathrm{N\,m}.
\]
第二姿态中重力为 $(-9.81,0,0)^{\mathsf T}$，则
\[
 f_2=(-19.62,0,0)^{\mathsf T},\qquad m_2=(0,-0.981,-0.981)^{\mathsf T}.
\]
若这些都是已去零偏、符号正确的静态载荷读数，可依次解出
\[
 m_p=-f_{1z}/9.81=2,\quad h_x=m_{1y}/9.81=0.2,
\]
\[
 h_y=-m_{1x}/9.81=-0.1,\quad h_z=-m_{2y}/9.81=0.1.
\]
再由 $c_p=h_p/m_p$ 得质心 $(0.1,-0.05,0.05)^{\mathsf T}\ \mathrm m$。第一种姿态单独不能看出沿重力方向的质心分量 $h_z$，第二种方向补足了信息。
\end{examplebox}

\begin{figure}[htbp]\centering
\includegraphics[width=0.90\linewidth]{figures/di_payload_gravity.pdf}
\caption{同一负载在不同法兰姿态下，重力在法兰坐标中的方向改变，从而暴露不同质心分量。图中力和力矩均按“负载作用于法兰”定义；几何为 $x$-$z$ 投影示意，箭头仅表示方向。}
\end{figure}

\subsection{多姿态数据怎样堆成线性方程}
将各姿态的 $A_p$ 竖向堆叠，得到 $w=A\beta_p$，其中 $\beta_p=(m_p,h_x,h_y,h_z)^{\mathsf T}$。使用上一章的 WLS 即可估计。若传感器还有未知固定零偏 $b_w$，应使用
\[
 w_i=A_{p,i}\beta_p+b_w+\epsilon_i,
\]
并把相应的常数列也加入回归。此时需要更多且方向丰富的数据，不能把一次静止读数的全部偏置都解释成质量。

使用关节电流推算负载时，同样可写 $\Delta\tau_i=-J_{F,i}^{\mathsf T}A_{p,i}\beta_p$。但负载前后差分还可能混入摩擦和线缆变化；静止关节中的静摩擦并不因两次数据相减就必然消失。可以比较正反接近同一姿态的结果，记录其差异，并通过多姿态冗余拟合检查残差。

需要多少姿态由矩阵的秩和条件数决定，不是固定的“三个”“六个”或“十个”。姿态很多却都绕同一轴小幅变化，仍可能缺少信息；两个方向在理想六维力传感器算例中足以解出四个组合，也不意味着真实系统两次测量就足够鲁棒。

\section{负载改变的力矩量级应先手算}
\begin{examplebox}[算例 10B 半公斤工具对两连杆重力的影响]
在本书两连杆末端再装质量 $0.5\ \mathrm{kg}$ 的工具，质心沿第二杆轴线超出末端 $0.1\ \mathrm m$。水平伸直时，质心到第一关节的力臂为 $0.5+0.5+0.1=1.1\ \mathrm m$，到第二关节为 $0.5+0.1=0.6\ \mathrm m$。
\[
 \Delta g_1=0.5\times9.81\times1.1=5.3955\ \mathrm{N\,m},
\]
\[
 \Delta g_2=0.5\times9.81\times0.6=2.9430\ \mathrm{N\,m}.
\]
原机械臂该姿态的重力补偿为 $(9.81,2.4525)$，加负载后为 $(15.2055,5.3955)\ \mathrm{N\,m}$。若 $q=(0,\pi/2)$，工具质心在第二关节正上方，第二关节附加重力力矩为零，第一关节则为 $0.5\times9.81\times0.5=2.4525\ \mathrm{N\,m}$。
\end{examplebox}

这个计算说明：同一个负载在不同姿态下造成的偏差不同。低阻抗下，近似平衡误差为模型残差除以刚度，因此“某个方向比较准，另一个方向回正不到位”可能来自姿态相关重力误差，而非单纯控制增益问题。

\begin{examplebox}[算例 10C 单位错误怎样放大负载力矩]
$1\ \mathrm{kg}$ 负载的水平力臂 $0.2\ \mathrm m$，重力矩为 $1.962\ \mathrm{N\,m}$。若误把 $20\ \mathrm{cm}$ 写成 $2\ \mathrm m$，就变成 $19.62\ \mathrm{N\,m}$。质量、距离、力矩三项应进行量纲和数量级检查，不能只相信程序返回了一个有限数。
\end{examplebox}

\section{把负载加入连杆之前先统一参考点}
已知负载在自身质心坐标系中的惯量 $I_{Cp}$，其质心和朝向在末连杆坐标系 $L$ 中分别为 $c_p^L$、$R_{LP}$。先旋转到连杆坐标，再平移到连杆原点：
\begin{align}
 I_{Cp}^L&=R_{LP}I_{Cp}R_{LP}^{\mathsf T},\\
 I_{Op}^L&=I_{Cp}^L+m_p\left(\|c_p^L\|^2I_3-c_p^L(c_p^L)^{\mathsf T}\right),\\
 h_p^L&=m_pc_p^L.
\end{align}
如果原末连杆标准参数已经关于同一原点、在同一坐标系，则质量、一阶矩和原点惯量可以逐项相加。若一个参数包关于法兰，另一个关于连杆原点，不能直接相加。组合后要写回质心惯量时，先算
\[
 m'=m_l+m_p,\qquad c'=\frac{h_l+h_p}{m'},
\]
再用反向平行轴定理
\[
 I_C'=I_O'-m'\left(\|c'\|^2I_3-c'(c')^{\mathsf T}\right).
\]
只关心重力时，质量和一阶矩足够；涉及快速转动和角加速度时，才需要可靠负载惯量。静态标定不应被包装为已完成全部十个负载参数辨识。

\section{关节零位与几何误差怎样线性化}
设几何参数向量 $\theta$ 包括关节零偏、连杆尺寸等，外部测得位置 $p^{\mathrm{meas}}$，模型为 $p=f(q,\theta)$。在当前参数 $\theta_0$ 附近作一阶展开：
\begin{equation}
 p^{\mathrm{meas}}-f(q,\theta_0)\approx
 J_\theta(q,\theta_0)\Delta\theta,\qquad
 J_\theta=\frac{\partial f}{\partial\theta}.
\end{equation}
把不同姿态的方程堆起来，又是最小二乘。更新参数后重新线性化，这就是常见的迭代非线性最小二乘思路。误差很大时一次线性修正可能不够，需要逐步更新而不能过度外推。

\subsection{零位误差就是雅可比的一种输入误差}
如果只估关节零偏 $\delta q$，局部有 $\Delta p\approx J(q)\delta q$。

\begin{examplebox}[算例 10D 由位置偏差估计两个关节零偏]
在本书标准姿态 $q=(0,\pi/2)$，测得相对模型的位置偏差 $\Delta p=(0.003,-0.002)^{\mathsf T}\ \mathrm m$。假设只存在关节零偏，则
\[
 \begin{bmatrix}-0.5&-0.5\\0.5&0\end{bmatrix}
 \begin{bmatrix}\delta q_1\\\delta q_2\end{bmatrix}
 =\begin{bmatrix}0.003\\-0.002\end{bmatrix}.
\]
第二行给 $\delta q_1=-0.004\ \mathrm{rad}$；第一行给 $-0.5(-0.004+\delta q_2)=0.003$，因此 $\delta q_2=-0.002\ \mathrm{rad}$。换成角度约为 $-0.2292^\circ$ 与 $-0.1146^\circ$。这个结果仅在“只有零偏”假设下成立；若工具长度和连杆长度也未知，一个姿态的两个位置分量无法区分全部误差来源。
\end{examplebox}

\subsection{长度误差与零偏可能互相补偿}
在同一标准姿态，模型末端位置为 $(l_1,l_2)$。如果只存在长度偏差，外测位置 $(0.502,0.497)\ \mathrm m$ 就对应 $\Delta l_1=2\ \mathrm{mm}$、$\Delta l_2=-3\ \mathrm{mm}$。但在另一个姿态，长度误差的作用方向会改变；多姿态数据才能检验这套解释。

还需选择“哪些参考坐标固定”。若同时允许基座坐标、末端工具坐标和部分连杆参数任意改变，可能出现多个参数组合产生相同测量结果的规范自由度。应固定参考基准，去掉不可辨识方向，而不是靠更多迭代消除数学歧义。

\begin{figure}[htbp]\centering
\includegraphics[width=0.90\linewidth]{figures/di_calibration_sensitivity.pdf}
\caption{零位偏差对不同两连杆姿态产生不同方向和大小的位置误差。右侧将误差放大显示，用多姿态数据才能区分一致的参数误差与偶发测量噪声。}
\end{figure}

\section{工具中心点标定与负载质心标定不是一回事}
TCP 是 tool center point，即工具中心点或实际工作点。它可以是针尖、夹爪接触点或相机光学中心。质心是质量分布的平均位置，两者一般不同。把 TCP 设准确可以改善位置命令，不能自动改善重力补偿。

一种常见 TCP 标定是让工具尖端始终抵住同一个固定点，改变法兰姿态。设法兰到工具尖端的未知位置为 $c_T$，固定点在基座中的位置为 $p_*$，每次法兰位姿为 $(R_i,t_i)$，则
\begin{equation}
 R_ic_T+t_i=p_*,\qquad
 \begin{bmatrix}R_i&-I_3\end{bmatrix}\begin{bmatrix}c_T\\p_*\end{bmatrix}=-t_i.
\end{equation}
堆叠多次测量即可求解。只绕一个轴转动时，该轴方向的信息可能不足；应使用足够多样的姿态，并检查残差是否来自尖端滑动或机械弹性。

\begin{examplebox}[算例 10E 二维固定点 TCP 标定]
为了便于手算，取二维旋转 $R_1=I$、$R_2=R(90^\circ)$，法兰位置 $t_1=(1,0)$、$t_2=(1.2,-0.2)\ \mathrm m$。两次工作点相同，所以
\[
 (I-R_2)c_T=t_2-t_1,\qquad
 \begin{bmatrix}1&1\\-1&1\end{bmatrix}c_T=\begin{bmatrix}0.2\\-0.2\end{bmatrix}.
\]
两行相加得 $2c_{Ty}=0$，因此 $c_{Ty}=0$；第一行给 $c_{Tx}=0.2$。固定工作点为 $p_*=t_1+c_T=(1.2,0)\ \mathrm m$。这是工具几何位置，不涉及工具质量。
\end{examplebox}

\section{手眼标定从坐标链推导到方程}
\subsection{先固定每个变换的方向}
对于眼在手上配置，相机刚性安装在法兰上。设
\[
 T_i={}^BT_{F_i},\qquad X={}^FT_C,\qquad C_i={} ^{C_i}T_{\mathrm{board}}.
\]
$T_i$ 来自机器人，$C_i$ 来自相机对固定标定板的位姿估计，$X$ 是要求的法兰到相机变换。固定标定板在基座中的位姿 $Z$ 应满足
\begin{equation}
 T_iXC_i=Z.
\end{equation}
对两次姿态 $i,j$ 有 $T_iXC_i=T_jXC_j$。左乘 $T_i^{-1}$、右乘 $C_j^{-1}$，整理得
\begin{equation}
 \underbrace{T_i^{-1}T_j}_{A_{ij}}X
 =X\underbrace{C_iC_j^{-1}}_{B_{ij}}.
\end{equation}
这就是 $AX=XB$。记住坐标链比背诵 $A$、$B$ 的文字定义更可靠，因为不同软件接口可能使用相反方向的输入。

\begin{figure}[htbp]\centering
\includegraphics[width=0.88\linewidth]{figures/di_handeye_chain.pdf}
\caption{眼在手上的坐标链。同一块静止标定板通过两次机器人与相机观测得到相同基座位姿，因此产生相对运动约束。}
\end{figure}

\subsection{先解旋转 再解平移}
把 $A,X,B$ 分别拆成旋转和平移，得到
\begin{align}
 R_AR_X&=R_XR_B,\\
 R_At_X+t_A&=R_Xt_B+t_X.
\end{align}
第一式等价于 $R_A=R_XR_BR_X^{\mathsf T}$。对不处于对数映射歧义的旋转，旋转向量满足 $a=R_Xb$，其中 $a=\operatorname{Log}(R_A)^\vee$、$b=\operatorname{Log}(R_B)^\vee$。对多组旋转向量，可求
\[
 \min_{R_X\in SO(3)}\sum_i\|a_i-R_Xb_i\|^2.
\]
令 $S=\sum_i a_ib_i^{\mathsf T}=U\Sigma V^{\mathsf T}$，则一种闭式旋转对齐解为
\[
 R_X=U\operatorname{diag}(1,1,\det(UV^{\mathsf T}))V^{\mathsf T}.
\]
对角修正确保行列式为正一，避免得到镜像反射。实际存在近 $180^\circ$ 运动或噪声时，还需处理旋转对数的分支和权重。

旋转确定后，平移方程是线性的：
\begin{equation}
 (R_A-I_3)t_X=R_Xt_B-t_A.
\end{equation}
把多组运动堆叠即可用最小二乘求 $t_X$。若全部运动只绕同一轴旋转，某些平移分量不可观察；若只有平移，$R_A-I=0$，$t_X$ 完全从方程中消失。丰富的平移可能帮助约束某些旋转信息，但不能由纯平移恢复完整手眼变换。

\begin{examplebox}[算例 10F 用具体矩阵验证 AX 与 XB]
取 $R_X=R_z(90^\circ)$，$t_X=(0.1,0.2,0.3)^{\mathsf T}\ \mathrm m$。设相机相对运动 $R_B=R_x(90^\circ)$、$t_B=(0.2,0,0)^{\mathsf T}$。由共轭关系得到 $R_A=R_y(90^\circ)$。再由平移关系得到
\[
 t_A=t_X+R_Xt_B-R_At_X=(-0.2,0.2,0.4)^{\mathsf T}.
\]
现在验证左边 $AX$ 的平移：
\[
 R_At_X+t_A=(0.3,0.2,-0.1)+(-0.2,0.2,0.4)=(0.1,0.4,0.3).
\]
右边 $XB$ 的平移为 $R_Xt_B+t_X=(0,0.2,0)+(0.1,0.2,0.3)=(0.1,0.4,0.3)$，完全相同。这个例子用于核对乘法方向；仅一组相对运动不足以唯一完成实际手眼标定。
\end{examplebox}

\section{相机内参与手眼误差如何影响机械臂定位}
针孔相机模型的核心为
\[
 u=f_x\frac XZ+c_x,\qquad v=f_y\frac YZ+c_y.
\]
$(X,Y,Z)$ 是点在相机系中的坐标，$f_x,f_y$ 以像素为单位表示焦距，$(c_x,c_y)$ 是主点，$(u,v)$ 是图像像素坐标。真实镜头还需畸变模型。相机内参不准确，标定板位姿 $C_i$ 就会带偏，随后污染手眼变换 $X$。

例如 $Z=1\ \mathrm m$、$f_x=600$ 像素，忽略深度误差时，横向像素误差 $2$ 像素对应 $\Delta X\approx Z\Delta u/f_x=0.003333\ \mathrm m$，约 $3.33\ \mathrm{mm}$。这只是局部针孔换算，不是完整三维定位误差预算；姿态、深度和标定板检测误差也会贡献。

验证手眼标定应使用没有参与求解的新姿态，检查 $T_iXC_i$ 是否在基座系中保持一致，并报告位置与旋转残差。位置残差与角度残差不能直接相加成一个没有说明尺度的数。小旋转误差还会通过工具长度放大为横向位置误差：$0.5^\circ$ 在 $0.3\ \mathrm m$ 长度上约产生 $0.3\times0.5\pi/180=2.62\ \mathrm{mm}$ 的位移。

\section{怎样把标定残差转成有用的误差预算}
标定后的一组参数不是终点，还需要知道它在哪些位置和姿态下可靠。对位置误差作局部线性化，可写
\[
 \delta p=J_q\delta q+J_\theta\delta\theta+J_T\delta\xi_T+\delta p_{\mathrm{meas}}.
\]
$\delta q$ 是关节测量或跟踪误差，$\delta\theta$ 是几何参数误差，$\delta\xi_T$ 是工具安装的小位姿误差，最后一项是外部测量误差。各个雅可比描述相应误差如何传到同一个位置坐标。所有量必须在相同参考系和相同时刻比较，否则“误差分解”本身又引入误差。

若各类误差可近似看成独立零均值随机量，位置协方差可由各项传播后相加。例如关节部分为 $\Sigma_p=J_q\Sigma_qJ_q^{\mathsf T}$。如果存在系统零偏或相关误差，就不能直接按平方和开根号缩小它们；有确定方向的偏置会持续偏向一侧。

这也解释了为何标定姿态应覆盖实际工作区。一个区域内数据丰富、参数组合可辨，另一个接近奇异或伸展极限的区域仍可能放大同样的参数误差。验证数据应包含实际使用的工具长度、负载方向和相机视角，而不只是与训练姿态相邻的一小团点。

对七轴相机云台，工具标定与手眼标定解决的是刚性安装关系。底座移动、支架弹性和镜头组件松动改变的是运行中的变换，不能靠一次静态标定永久消除。若误差随着姿态、负载或时间系统变化，先判断“固定参数有误”还是“参数确实在变化”，再选择标定、在线观测或机械加固。

\section{计算练习与分步解答}
\subsection*{练习 10.1 从一阶矩恢复质心}
估得质量 $1.5\ \mathrm{kg}$、一阶矩 $(0.12,-0.03,0.075)^{\mathsf T}\ \mathrm{kg\,m}$，求质心。

\textbf{解答。}逐分量除以质量，$c=(0.08,-0.02,0.05)^{\mathsf T}\ \mathrm m$。若质量估计误差明显，三个质心分量都会受同一个分母影响，报告不确定性时不能把它们当作独立误差。

\subsection*{练习 10.2 转换重力方向}
基座重力为 $(0,0,-9.81)$，法兰相对基座绕 $y$ 轴旋转 $90^\circ$。求法兰系重力。

\textbf{解答。}$R_{BF}=\begin{bmatrix}0&0&1\\0&1&0\\-1&0&0\end{bmatrix}$，所以
$R_{BF}^{\mathsf T}(0,0,-9.81)^{\mathsf T}=(9.81,0,0)^{\mathsf T}$。这里应乘转置；如果乘 $R_{BF}$，会得到相反的 $x$ 方向。

\subsection*{练习 10.3 工具误差与质心误差分别影响什么}
一公斤工具的质心水平距离模型少算 $0.03\ \mathrm m$，同时 TCP 的横向位置少算 $0.005\ \mathrm m$。分别计算重力矩偏差和工具几何位置偏差。

\textbf{解答。}质心误差造成重力矩偏差 $1\times9.81\times0.03=0.2943\ \mathrm{N\,m}$。TCP 几何位置偏差为 $5\ \mathrm{mm}$。修正 TCP 不会自动消除 $0.2943\ \mathrm{N\,m}$ 的重力偏差；修正质心也不自动把工作点移动到正确位置。

\subsection*{练习 10.4 纯平移为什么不能标定相机安装平移}
在 $AX=XB$ 中令 $R_A=R_B=I$，写出剩余约束。

\textbf{解答。}旋转方程变成 $R_X=R_X$，平移方程变为 $t_X+t_A=R_Xt_B+t_X$，消去 $t_X$ 后得到 $t_A=R_Xt_B$。因此相机安装平移 $t_X$ 没有出现在数据关系中，无论再采多少这种纯平移，仍不能通过它确定 $t_X$。

\section*{本章与原稿的对应}
本章扩展 [B，第44--45页、第46--50页] 的负载、运动学与手眼标定，并与 [A，第22--30页] 的坐标变换及运动学对应。重点补充了重力力旋量的符号、静态四参数回归、负载参数同坐标合并、TCP 与质心区别，以及 $AX=XB$ 的坐标链和数值验算。
